%%%%%%%%%%%%%%%%%%%%%%%%%%%%%%%%%%%%%%%%%
% Medium Length Professional CV
% LaTeX Template
% Version 3.0 (December 17, 2022)
%
% This template originates from:
% https://www.LaTeXTemplates.com
%
% Author:
% Vel (vel@latextemplates.com)
%
% Original author:
% Trey Hunner (http://www.treyhunner.com/)
%
% License:
% CC BY-NC-SA 4.0 (https://creativecommons.org/licenses/by-nc-sa/4.0/)
%
%%%%%%%%%%%%%%%%%%%%%%%%%%%%%%%%%%%%%%%%%

%----------------------------------------------------------------------------------------
%	PACKAGES AND OTHER DOCUMENT CONFIGURATIONS
%----------------------------------------------------------------------------------------

\documentclass[
	%a4paper, % Uncomment for A4 paper size (default is US letter)
	11pt, % Default font size, can use 10pt, 11pt or 12pt
]{resume} % Use the resume class
\usepackage{fontawesome5}
\usepackage{hyperref}
\usepackage{marvosym}
\usepackage{multicol}

\hypersetup{
    colorlinks=true,
    linkcolor=black,   % Color for normal internal links.
    urlcolor=black,    % Color for URLs.
    citecolor=black    % Color for citation links.
}

% \usepackage{ebgaramond} % Use the EB Garamond font

%------------------------------------------------

\name{Yanqi Su} % Your name to appear at the top

% You can use the \address command up to 3 times for 3 different addresses or pieces of contact information
% Any new lines (\\) you use in the \address commands will be converted to symbols, so each address will appear as a single line.

% \address{School of Computing, The Australian National University} 
\address{School of Computation, Information and Technology, Technical University of Munich} 
% Main address

% \address{123 Pleasant Lane \\ City, State 12345} % A secondary address (optional)

\address{yanqi.su@tum.de \\ (+49) 151 58748540} % Contact information

%----------------------------------------------------------------------------------------

\begin{document}

%----------------------------------------------------------------------------------------
%	EDUCATION SECTION
%----------------------------------------------------------------------------------------

\begin{rSection}{Research Interests}
	
% My research interests are \textbf{Software Testing}, \textbf{Program Analysis}, \textbf{Knowledge Graph} on Software Engineering, \textbf{Software Repository Mining}, \textbf{Multi-Modal Models}, \textbf{Human-Computer Interaction}.


My research interests include \textbf{Software Testing}, \textbf{Program Analysis}, \textbf{Knowledge Graphs} and \textbf{Large Language Models (LLMs)} for Software Engineering, and  \textbf{Software Repository Mining}.
% , and \textbf{Human-Computer Interaction}.
	
\end{rSection}

\begin{rSection}{Education}
	
	\textbf{The Australian National University (ANU)} \hfill \textit{Sep. 2020 - Sep. 2024} \\ 
	Ph.D. in Computer Science\\
    Supervisor: Zhenchang Xing \\\\
    \textbf{Nanjing University of Aeronautics and Astronautics (NUAA)} \hfill \textit{Sep. 2017 - Jun. 2020} \\ 
	M.E. in Computer Science\\
     % Supervisor: Yu Zhou \\
     \\
 % \& Engineering \\
    \textbf{Nanjing University of Aeronautics and Astronautics (NUAA)} \hfill \textit{Sep. 2013 - Jun. 2017} \\ 
	B.E. in Computer Science
 
	% Minor in Linguistics \smallskip \\
	% Member of Eta Kappa Nu \\
	% Member of Upsilon Pi Epsilon \\
	% Overall GPA: 5.678
	
\end{rSection}

\begin{rSection}{Publications}
 \href{https://arxiv.org/pdf/2603.03121}{\textbf{Yanqi Su}, Michael Pradel, Chunyang Chen: RippleGUItester: Change-Aware Exploratory Testing. The 35th SIGSOFT International Symposium on Software Testing and Analysis (ISSTA 2026).} Acceptance rate: 10.1\% (90/888)
\newline\newline
 \href{https://dl.acm.org/doi/10.1145/3715752}{\textbf{Yanqi Su}, Zhenchang Xing, Chong Wang, Chunyang Chen, Xiwei Xu, Qinghua Lu, Liming Zhu: Automated Soap Opera Testing Directed by LLMs and Scenario Knowledge: Feasibility, Challenges, and Road Ahead. The ACM International Conference on the Foundations of Software Engineering (FSE2025). }
\newline\newline
 \href{https://conf.researchr.org/details/icse-2024/icse-2024-research-track/157/Enhancing-Exploratory-Testing-by-Large-Language-Model-and-Knowledge-Graph}{\textbf{Yanqi Su}, Dianshu Liao, Zhenchang Xing, Qing Huang, Mulong Xie, Qinghua Lu, Xiwei Xu:
Enhancing Exploratory Testing by Large Language Model and Knowledge Graph. The 46th International Conference on Software Engineering (ICSE'24). }
\newline\newline
\href{https://ieeexplore.ieee.org/document/10172537}{\textbf{Yanqi Su}, Zheming Han, Zhenchang Xing, Xiwei Xu, Liming Zhu, Qinghua Lu:
SoapOperaTG: A Tool for System Knowledge Graph Based Soap Opera Test Generation. The 45th International Conference on Software Engineering (ICSE'23 Tool Demo).}
\newline\newline
\href{https://dl.acm.org/doi/abs/10.1145/3551349.3556967}{\textbf{Yanqi Su}, Zheming Han, Zhenchang Xing, Xin Xia, Xiwei Xu, Liming Zhu, Qinghua Lu:
Constructing a System Knowledge Graph of User Tasks and Failures from Bug Reports to Support Soap Opera Testing. 37th IEEE/ACM International Conference on Automated Software Engineering (ASE2022).}\newline\newline
\href{https://ieeexplore.ieee.org/document/9678574}{\textbf{Yanqi Su}, Zhenchang Xing, Xin Peng, Xin Xia, Chong Wang, Xiwei Xu, Liming Zhu:
Reducing Bug Triaging Confusion by Learning from Mistakes with a Bug Tossing Knowledge Graph. 36th IEEE/ACM International Conference on Automated Software Engineering (ASE2021). \newline
\faMedal \hspace{0.5em}\textbf{Distinguished Paper Award}}\newline\newline
\href{https://ieeexplore.ieee.org/document/10174001}{\textbf{Yanqi Su}, Zheming Han, Zhipeng Gao, Zhenchang Xing, Qinghua Lu, Xiwei Xu:
Still Confusing for Bug-Component Triaging? Deep Feature Learning and Ensemble Setting to Rescue. The 31st IEEE/ACM International Conference on Program Comprehension (ICPC2023).} 
\newline\newline
% \href{https://ieeexplore.ieee.org/document/8961125}{Yu Zhou, \textbf{Yanqi Su}, Taolue Chen, Zhiqiu Huang, Harald C. Gall, Sebastiano Panichella:
% User Review-Based Change File Localization for Mobile Applications. IEEE Trans. on Software Engineering. Vol.47, No.12, pp:2755-2770, IEEE, 2021. (TSE)}\newline\newline
\href{https://ieeexplore.ieee.org/document/8961125}{Yu Zhou, \textbf{Yanqi Su}, Taolue Chen, Zhiqiu Huang, Harald C. Gall, Sebastiano Panichella:
User Review-Based Change File Localization for Mobile Applications. IEEE Trans. on Software Engineering. (TSE)}\newline\newline
Zhipeng Gao, \textbf{Yanqi Su}, Xing Hu, Xin Xia: \href{https://dl.acm.org/doi/10.1145/3700793}{Automating TODO-dropped Methods Detection and Patching. ACM Transactions on Software Engineering and Methodology. (TOSEM)} 
% \newline\newline
% \textbf{Yanqi Su}, Dianshu Liao, Zhenchang Xing, Qing Huang, Mulong Xie, Qinghua Lu, Xiwei Xu:
% Enhancing Exploratory Testing by Large Language Model and Knowledge Graph. The 46th International Conference on Software Engineering (ICSE'24). 
\end{rSection}

\begin{rSection}{Teaching and Working}

\textbf{Postdoctoral Researcher in Technische Universität München (TUM)}, Software Engineering \& AI, from Oct-21-2024, supervised by Chunyang Chen
\newline\newline
\textbf{Lecturer in TUM}, Seminar Exploratory Software Testing (INHN0015, INHN4016, IN0014, IN45111), 2026 Summer
\newline\newline
\textbf{Lecturer in TUM}, Seminar Exploratory Software Testing (IN0014) (INHN0015, INHN4012), 2025 Summer
\newline\newline
\textbf{Guest Lecturer in TUM}, Advanced Topics in Software Engineering (CITHN3003), 2024 \& 2025 Winter
\newline\newline
\textbf{Teaching Assistant in ANU}, COMP4130 Managing Software Quality and Process, 2024S1, supervised by Xiaoyu Sun
\newline\newline
% (Email: Xiaoyu.sun1@anu.edu.au)}
%\textbf{Internship in Huawei}, Software Engineering Application Technology Lab, from Jun-30-2021 to Sep-30-2022, supervised by Xin Xia
%\newline\newline
% (Email: Xin.Xia@acm.org)}
\textbf{Research Visit in University of Stuttgart \& CISPA}, Software Lab, from Nov-01-2025 to Nov-12-2025, supervised by Michael Pradel
\newline\newline
\textbf{Research Assistant in Fudan University}, CodeWisdom Team, from Sep-01-2020 to Jun-30-2021, supervised by Xin Peng
\newline\newline
% (Email: pengxin@fudan.edu.cn)}
\textbf{Teaching Assistant in NUAA}, Software Metrics, 2017S2, supervised by Yu Zhou

\end{rSection}

\begin{rSection}{Mentor}
    \begin{description}
      \item[2025 \textbf{Abdelrahman Diab}, TUM] Master thesis: Bug-Inducing Code Change Detection
      \item[2025 \textbf{Lazar Minic}, TUM] Master thesis: Bridging Policy and Practice: Leveraging LLMs for Privacy-Aware Exploratory Software Testing
      \item[2024 \textbf{Alper Temel}, TUM] Master thesis: Non-intrusive GUI Element Location Using LLMs
      \item[2024 \textbf{Linsen Mao}, TUM] Bachelor thesis: Soap Opera Testing on Enterprise Resource Planning Software
      \item[2024 \textbf{Bryan Alvin}, TUM] Bachelor thesis: Automated Exploratory Testing of Business Logic in ERP Systems
      \item[2022 \textbf{Zheming Han}, ANU] Master thesis: Reducing Bug Component Triaging Mistakes by Deep Learning
      \item[2017 \textbf{Jian Liu}, NUAA] Bachelor thesis: Code Reviewer Recommendation Based on File Path Analysis
    \end{description}
    % 2025 Lazar Minic, TUM, Master thesis: Bridging Policy and Practice: Leveraging LLMs for Privacy-Aware Exploratory Software Testing\newline\newline
    % 2024 Alper Temel, TUM, Master thesis: Non-intrusive GUI Element Location Using LLMs\newline\newline
    % 2024 Linsen Mao, TUM, Bachelor thesis: Soap Opera Testing on Enterprise Resource Planning Software\newline\newline
    % 2024 Bryan Alvin, TUM, Bachelor thesis: Automated Exploratory Testing of Business Logic in ERP Systems\newline\newline
    % 2022 Zheming Han, ANU, Master thesis: Reducing Bug Component Triaging Mistakes by Deep Learning
    % \newline\newline
    % 2017 Jian Liu, NUAA, Bachelor thesis: Code Reviewer Recommendation Based on File Path Analysis
    
\end{rSection}

\begin{rSection}{Fundings \& Awards}
ANU Vice-Chancellor's Travel Grant
\newline\newline
ACM SIGSOFT CAPS AWARD FSE/ISSTA 2025 
\newline\newline
ACM Travel Grant for 2024 IEEE/ACM 46th International Conference on Software Engineering
\newline\newline
ACM SIGSOFT Distinguished Paper Award (First author) in ASE2021
\newline\newline
HDR Fee Remission Merit Scholarship, ANU, 2020.09 – 2024.09
\newline\newline
University Research Scholarship, ANU, 2020.09 – 2024.09
\newline\newline
First-Class Graduate Scholarship, NUAA, 2017.09 – 2020.04
\newline\newline
Outstanding Research and Innovation Award, NUAA, 2019.09 – 2020.04
\newline\newline
Outstanding Contribution to Social Activities, NUAA, 2019.09 – 2020.04
\newline\newline
Third-Class Scholarship, NUAA, 2015.09 – 2017.06
\newline\newline
Second-Class Scholarship, NUAA, 2014.09 – 2015.06
\newline\newline
Third-Class Scholarship, NUAA, 2013.09 – 2014.06
\newline\newline
Microsoft Research Asia Club Star, 2014
\newline\newline
Outstanding Youth Volunteer, 2nd Youth Olympic Games, 2014.08 – 2014.09
\end{rSection}

\begin{rSection}{Service}
\textbf{Journal Reviewer}
\begin{itemize}
    \item IEEE Transactions on Software Engineering (TSE)
    \item ACM Transactions on Software Engineering and Methodology (TOSEM)
    \item Empirical Software Engineering (EMSE)
    \item Software Testing, Verification and Reliability (STVR)
\end{itemize}
\textbf{Conference Program Committee}
\begin{itemize}
    \item The ACM International Conference on the Foundations of Software Engineering (FSE2027)
    \item The 49th IEEE/ACM International Conference on Software Engineering (ICSE2027)
    \item The 41th IEEE/ACM International Conference on Automated Software Engineering (ASE2026)
    \item The 23rd International Conference on Mining Software Repositories (MSR2026)
    \item The International Conference on Predictive Models and Data Analytics in Software Engineering (PROMISE 2026) 
    \item The 1st International Workshop on User Interface and Experience for Software Engineering (UISE 2026) 
    \item The 40th IEEE/ACM International Conference on Automated Software Engineering (ASE2025)
    \item The 36th IEEE International Symposium on Software Reliability Engineering (ISSRE2025)
    \item The 32nd Asia-Pacific Software Engineering Conference (APSEC2025 SEIP Track \& Student Research Competition Track)
    % \item The 32nd Asia-Pacific Software Engineering Conference (APSEC 2025 Student Research Competition Track)
\end{itemize}
\textbf{Student Volunteer}
\begin{itemize}
    \item The ACM International Conference on the Foundations of Software Engineering (FSE2025)
\end{itemize}
% \textbf{Journal Reviewer} ACM Transactions on Software Engineering and Methodology (TOSEM)
% \newline\newline
% \textbf{Journal Reviewer} Empirical Software Engineering (EMSE)
% \newline\newline
% \textbf{Program Committee} 
% \newline\newline
% \textbf{Program Committee} 40th IEEE/ACM International Conference on Automated Software Engineering (ASE2025)
% \newline\newline
% \textbf{Program Committee} 36th IEEE International Symposium on Software Reliability Engineering (ISSRE2025)
% \newline\newline
% \textbf{Student Volunteer} ACM International Conference on the Foundations of Software Engineering (FSE2025) 
\end{rSection}

\begin{rSection}{Talk}
\textbf{Shaping the Future of Exploratory Software Testing with Knowledge Graphs and LLMs}
\newline
This talk is a compilation of my main research and I gave this talk at multiple places.
\begin{itemize}
    % \item Jul. 2025 Huazhong University of Science and Technology, China
    \item Jul. 2025 The Chinese University of Hong Kong, Shenzhen, China
    \item Jun. 2025 University of Oxford, UK
    \item Apr. 2025 Aalto University, Finland
\end{itemize}
\textbf{Automated Soap Opera Testing Directed by LLMs and Scenario Knowledge}
\begin{itemize}
    \item Feb. 2025 Technical University of Munich, Germany
\end{itemize}
\textbf{Constructing a System Knowledge Graph of User Tasks and Failures to Support Soap Opera Testing}
\begin{itemize}
    \item Dec. 2023 Google Research Australia, Australia
    % \newline
\end{itemize}


% \textbf{Huazhong University of Science and Technology} Shaping the Future of Exploratory Software Testing with Knowledge Graphs and LLMs, Jul. 2025
% \newline\newline
% \textbf{The Chinese University of Hong Kong, Shenzhen} Shaping the Future of Exploratory Software Testing with Knowledge Graphs and LLMs, Jul. 2025
% \newline\newline
% \textbf{University of Oxford} Shaping the Future of Exploratory Software Testing with Knowledge Graphs and LLMs, Jun. 2025
% \newline\newline
% \textbf{Aalto University} Shaping the Future of Exploratory Testing with Knowledge Graphs and LLMs: Comprehensive, Adaptive, and Intelligent, Apr. 2025
% \newline\newline
% \textbf{Technical University of Munich} Automated Soap Opera Testing Directed by LLMs and Scenario Knowledge, Feb. 2025
% \textbf{Google Research Australia} Constructing a System Knowledge Graph of User Tasks and Failures to Support Soap Opera Testing, Dec. 2023
\end{rSection}

\begin{rSection}{References}

\begin{multicols}{2}

\textbf{Zhenchang Xing} (Supervisor)\\
Associate Professor\\
School of Computing\\
The Australian National University\\
Senior Principal Research Scientist\\ 
CSIRO's Data61\\
Australia\\[0.2em]
Zhenchang.Xing@data61.csiro.au\\[1em]

\textbf{Chunyang Chen}\\
Professor\\
School of Computation, Information and Technology\\
Technical University of Munich, Germany\\[0.2em]
chun-yang.chen@tum.de

\columnbreak

\textbf{Xin Peng}\\
Deputy Dean \& Professor\\
School of Computer Science\\
Fudan University, China\\[0.2em]
pengxin@fudan.edu.cn

\end{multicols}
\end{rSection}


% \textbf{Zhenchang Xing}, Associate Professor, School of Computing, Australian National University, Australia

% \Letter{} Zhenchang.Xing@data61.csiro.au
% \newline\newline
% % \textbf{Xiaoyu Sun}, Lecturer, School of Computing, Australian National University, Australia
% % \Letter{} Xiaoyu.Sun1@anu.edu.au
% \textbf{Chunyang Chen}, Professor, School of Computation, Information and Technology, Technical University of Munich, Germany

% \Letter{} chun-yang.chen@tum.de
% \newline\newline
% \textbf{Xin Peng}, Deputy Dean \& Professor, School of Computer Science, Fudan University, China

% \Letter{} pengxin@fudan.edu.cn
% \newline\newline
% \textbf{Qinghua Lu}, Principal Research Scientist, Responsible AI Science Team, CSIRO's Data61, Australia

% \Letter{} qinghua.lu@data61.csiro.au




%\newline\newline
%\textbf{Xin Xia}, Director of the Software Engineering Application Technology Lab, Huawei, China

%\Letter{} Xin.Xia@acm.org

    

%----------------------------------------------------------------------------------------
%	WORK EXPERIENCE SECTION
%----------------------------------------------------------------------------------------

% \begin{rSection}{Experience}

% 	\begin{rSubsection}{ACME, Inc}{October 2010 - Present}{Web Developer}{Palo Alto, CA}
% 		\item Lorem ipsum dolor sit amet, consectetur adipiscing elit. Donec a diam lectus.
% 		\item Donec et mollis dolor. Praesent et diam eget libero Adobe Coldfusion egestas mattis sit amet vitae augue.
% 		\item Nam tincidunt congue enim, ut porta lorem Microsoft SQL lacinia consectetur.
% 		\item Donec ut libero sed arcu vehicula ultricies a non tortor. Lorem ipsum dolor sit amet, consectetur adipiscing elit.
% 		\item Pellentesque auctor nisi id magna consequat JavaScript sagittis.
% 		\item Aliquam at massa ipsum. Quisque bash bibendum purus convallis nulla ultrices ultricies.
% 	\end{rSubsection}

% %------------------------------------------------

% 	\begin{rSubsection}{AJAX Hosting}{December 2009 - October 2010}{Lead Developer}{Austin, TX}
% 		\item Aenean ut gravida lorem. Ut turpis felis, Perl pulvinar a semper sed, adipiscing id dolor.
% 		\item Curabitur dapibus enim sit amet elit pharetra tincidunt website feugiat nisl imperdiet. Ut convallis AJAX libero in urna ultrices accumsan.
% 		\item Cum sociis natoque penatibus et magnis dis MySQL parturient montes, nascetur ridiculus mus.
% 		\item In rutrum accumsan ultricies. Mauris vitae nisi at sem facilisis semper ac in est.
% 		\item Nullam cursus suscipit nisi, et ultrices justo sodales nec. Fusce venenatis facilisis lectus ac semper.
% 	\end{rSubsection}

% %------------------------------------------------

% 	\begin{rSubsection}{TinySoft}{January 2008 - April 2010}{Web Designer \& Developer}{Gainesville, GA}
% 		\item Vivamus PostgreSQL fermentum semper porta. Nunc diam velit PHP, adipiscing ut tristique vitae
% 		\item Maecenas convallis ullamcorper ultricies stylesheets.
% 		\item Quisque mi metus, unit tests CSS ornare sit amet fermentum et, tincidunt et orci.
% 		\item Curabitur venenatis pulvinar tellus gravida ornare. Sed et erat faucibus nunc euismod ultricies ut id
% 	\end{rSubsection}

% \end{rSection}

% %----------------------------------------------------------------------------------------
% %	TECHNICAL STRENGTHS SECTION
% %----------------------------------------------------------------------------------------

% \begin{rSection}{Technical Strengths}

% 	\begin{tabular}{@{} >{\bfseries}l @{\hspace{6ex}} l @{}}
% 		Computer Languages & Prolog, Haskell, AWK, Erlang, Scheme, ML \\
% 		Protocols \& APIs & XML, JSON, SOAP, REST \\
% 		Databases & MySQL, PostgreSQL, Microsoft SQL \\
% 		Tools & SVN, Vim, Emacs
% 	\end{tabular}

% \end{rSection}

% %----------------------------------------------------------------------------------------
% %	EXAMPLE SECTION
% %----------------------------------------------------------------------------------------

% %\begin{rSection}{Section Name}

% 	%Section content\ldots

% %\end{rSection}

% %----------------------------------------------------------------------------------------

\end{document}
